\documentclass[12pt]{article}
\usepackage[utf8]{inputenc}
\usepackage{amsmath,amssymb,amsthm}
\usepackage[margin=1in]{geometry}

\newtheorem{theorem}{Theorem}
\newtheorem{lemma}[theorem]{Lemma}
\newtheorem{corollary}[theorem]{Corollary}
\newtheorem{proposition}[theorem]{Proposition}
\newtheorem{definition}[theorem]{Definition}

\title{Problem 49: Prime Permutations}
\author{Project Euler Solution}
\date{}

\begin{document}
\maketitle

\section{Problem Statement}

The arithmetic sequence $1487, 4817, 8147$ (common difference $3330$) consists of three 4-digit primes that are digit permutations of one another. There is one other such sequence of three 4-digit primes. What 12-digit number is formed by concatenating its three terms?

\section{Formal Development}

\begin{definition}[Digit Permutation Equivalence]
For non-negative integers $a, b$, write $a \sim b$ if they share the same multiset of decimal digits. The \emph{canonical signature} $\sigma(n)$ is the string formed by sorting the digits of $n$ in non-decreasing order. By definition, $a \sim b \iff \sigma(a) = \sigma(b)$.
\end{definition}

\begin{definition}[AP-Permutation Triple]
A triple $(p_1, p_2, p_3)$ with $p_1 < p_2 < p_3$ is an \emph{AP-permutation triple} if all three are prime, $p_2 - p_1 = p_3 - p_2$, and $p_1 \sim p_2 \sim p_3$.
\end{definition}

\begin{lemma}[Common Difference Bound]\label{lem:dbound}
For a 4-digit AP-permutation triple, the common difference $d$ satisfies $d \leq 4499$.
\end{lemma}

\begin{proof}
We need $1000 \leq p_1$ and $p_1 + 2d \leq 9999$, so $2d \leq 8999$.
\end{proof}

\begin{lemma}[Grouping]\label{lem:group}
All members of an AP-permutation triple share the same canonical signature $\sigma$.
\end{lemma}

\begin{proof}
Immediate from the digit permutation condition.
\end{proof}

\begin{theorem}[Complete Enumeration]\label{thm:main}
The only 4-digit AP-permutation triples are:
\begin{enumerate}
\item $(1487, 4817, 8147)$ with $d = 3330$, and
\item $(2969, 6299, 9629)$ with $d = 3330$.
\end{enumerate}
\end{theorem}

\begin{proof}
Generate all $\pi_4 = 1061$ four-digit primes by sieve. Partition them into equivalence classes under $\sim$ via $\sigma$. For each class $G$ with $|G| \geq 3$, examine all pairs $(p_i, p_j)$ with $p_i < p_j$ and test whether $2p_j - p_i \in G$. The exhaustive search yields exactly the two triples listed.

Verification for the second triple:
\begin{itemize}
\item $2969$, $6299$, $9629$ are each prime (by sieve or trial division).
\item $6299 - 2969 = 9629 - 6299 = 3330$.
\item $\sigma(2969) = \sigma(6299) = \sigma(9629) = \text{``2699''}$.
\end{itemize}
No other triple exists, as the search is exhaustive over all classes.
\end{proof}

\begin{corollary}
The answer is the concatenation $296962999629$.
\end{corollary}

\section{Complexity Analysis}

\begin{proposition}[Time]
$O(N \log \log N + \pi_4 \cdot g_{\max}^2)$, where $N = 9999$, $\pi_4 = 1061$, and $g_{\max}$ is the maximum class size. In practice, $g_{\max} \leq 15$, so the second term is $O(\pi_4)$.
\end{proposition}

\begin{proposition}[Space]
$O(\pi_4)$ for the prime groups and set lookup structures.
\end{proposition}

\section{Answer}

\[
\boxed{296962999629}
\]

\end{document}
